\section{Regras e Inferência}
\label{sec:regras}

A aplicação de regras é peça fundamental para o funcionamento do eProfile. Uma regra é uma instrução de inferência registrada por uma aplicação que será executada sobre trilhas de uma entidade, resultando num perfil de entidade de acordo com um modelo de entidade previamente cadastrado pela aplicação.

Para que uma regra possa ser aplicada, as seguintes informações precisam ter sido previamente registradas no eProfile:

\begin{itemize}

\item a regra propriamente dita;

\item um modelo de entidade, que servirá como base para o formato de perfil da entidade;

\item a trilha da entidade, sobre a qual a inferência será executada.

\end{itemize}

A partir destas três informações, o raciocinador tem condições de gerar um perfil de entidade, de acordo com a Figura~\ref{fig:dados_inferencia}.

\begin{figure}
  \caption{Fluxo de dados para inferência do perfil de entidade}
  \includegraphics[width=0.8\textwidth]{imagens/dados_inferencia.png}
  \label{fig:dados_inferencia}
  
Fonte: elaborado pelo autor
\end{figure}

As regras são sempre definidas pela aplicação, e uma regra é composta pelo seguinte conjunto de informações:

\begin{itemize}

\item {\bf Nome:} um identificador único para a regra, que permite que a aplicação identifique a regra posteriormente, para que possa editá-la ou excluí-la;

\item {\bf Campo do perfil:} define qual é o campo do perfil da entidade que deverá ser atualizado através da execução desta inferência. Este campo deve ser obrigatoriamente equivalente a uma classe definida na ontologia de modelo de entidade;

\item {\bf Regra:} a regra que será executada sobre a trilha da entidade;

\item {\bf Atualização:} uma instrução que define como o campo será atualizado (por exemplo, o valor pode ser simplesmente sobrescrito ou atualizado de acordo com alguma fórmula);

\item {\bf Frequência de atualização:} define a periodicidade de execução da regra. É importante notar que o valor definido aqui pode não ser honrado, devido a outros motivos (como uma sobrecarga no servidor);

\item {\bf Validade}: define qual a data limite para a execução desta regra. A ausência desta informação define que a regra continuará sendo processada indefinidamente.

\end{itemize}

O raciocinador gerenciará a carga de processamento no servidor. Assim, algumas regras podem não ser executadas de acordo com o agendamento solicitado. Além disso, é possível que uma aplicação registre certas regras mas deixe de usar o servidor, fazendo com que as regras continuem sendo executadas desnecessariamente. Assim, caso um perfil não seja consultado por um tempo determinado, a frequência da execução das regras passa a diminuir progressivamente, até o ponto em que a regra não seja mais executada. Caso a aplicação volte a consultar os perfis, a frequência de execução das regras é normalizada.

As regras podem ser escritas em duas linguagens diferentes: SPARQL ou SWRL. O SPARQL \cite{sparql} é uma linguagem de consulta de ontologias cuja sintaxe assemelha-se a linguagem de consulta de banco de dados SQL \cite{sql}, enquanto o SWRL \cite{swrl} é uma linguagem de regras que permite a consulta utilizando construtores lógicos mais complexos que o SPARQL. Além disso, as informações podem ser inferidas através das relações axiomáticas entre as classes.

Uma vez que a regra tenha sido validada pelo agente configurador, ela entra em funcionamento e passa a ser elegível para execução.

Tanto as consultas via SPARQL/SWRL quanto a geração de classes a partir das relações axiomáticas são executadas através de um motor de inferências externo que é descrito na Seção~\ref{sec:motor_inferencia}. Um exemplo detalhado da aplicação de regras pode ser verificado na Seção~\ref{sec:exemplo}.
